\section{Discussion}
\label{sec:discussion}

\paragraph{Implications for deployment.}
A uniform compression budget implicitly accepts order-of-magnitude severity penalties during low-sparsity reasoning phases. Pilot evaluation ($n{=}30$, $s{=}0.25$; Appendix~\ref{app:pilot-mitigation}) achieved 18\% oracle recovery (95\% CI: $[-39\%, 45.8\%]$), falling below the pre-registered 30\% gate; the CI spans zero, indicating no statistically distinguishable improvement. Effective phase-aware allocation likely requires finer phase boundary estimation than current proxies provide. Additional mechanisms (e.g., strategy switching, adaptive thresholding) warrant investigation. We hypothesize that phase-modulated severity may extend to alternative attention-mass-based scoring methods (H$_2$O \cite{h2o-neurips2023}, SnapKV \cite{snapkv-neurips2024}, PyramidKV \cite{pyramidkv-colm2025}), which remains an open question for future work. Extended discussion in Appendix~\ref{app:extended-discussion}.

\paragraph{Reasoning depth scaling and architecture dependence.}
Strategy specificity is not a fixed property but varies along at least two axes. First, it scales with reasoning chain depth: MATH-500 ($15$--$17{\times}$) $>$ GSM8K (${\sim}2{\times}$ point-estimate ratio; random non-significant at $n{=}50$), with within-MATH-500 correlational support (Spearman $\rho > 0$, $p < 0.02$ between trace length and specificity). Second, it is architecture-dependent: LlamaForCausalLM exhibits significant random phase gradients ($d_z{=}0.41$--$0.46$), reducing the importance/random ratio to $2.0$--$2.6{\times}$ compared with $15$--$17{\times}$ on DeepSeek (Qwen2-based). This suggests that the degree to which random eviction is phase-insensitive depends on architectural details of the attention implementation. The depth-scaling pattern is consistent with the discrimination collapse mechanism: longer low-sparsity phases provide more tokens over which pruning-degraded importance scores accumulate errors, amplifying the gap between importance-based and random eviction. The practical implication holds regardless of architecture: as models generate longer reasoning chains (a trend driven by test-time compute scaling~\cite{snell-2024-scaling-test-time}), the penalty for naive compound compression grows superlinearly, though the relative advantage of strategy-aware allocation varies by model family. Behavioral evidence supports this scaling concern: under compound sparsity at $s \geq 0.25$, DeepSeek-R1-Distill-Qwen-7B generates degenerate reasoning loops averaging 6{,}344 tokens (76\% exceeding the 8{,}192-token budget), suggesting that compound sparsity not only amplifies phase-internal severity but may also disrupt reasoning termination. This degeneration is consistent with the discrimination collapse observed in exploration phases: once importance scores lose discriminative power, the model fails to consolidate toward an answer, producing unbounded exploratory chains.

\paragraph{Diagnostic framework value.}
While we do not propose a new compression method, the phase-gated analysis framework has direct engineering utility: (1)~the Interaction Index can serve as a per-phase evaluation metric for assessing the phase robustness of any new compound compression scheme; (2)~the dose-response curves (Figure~\ref{fig:hero}) provide target operating points for phase-aware controllers, e.g., maintaining $s \leq 0.20$ during low-sparsity phases to keep $\Delta < 1.5$; (3)~the strategy specificity finding carries a concrete engineering implication: importance-based eviction under compound sparsity requires phase-aware protection that random eviction does not, because the interaction is structurally coupled to the scoring mechanism rather than to compression magnitude alone.

\paragraph{Recency control as methodological contribution.}
The use of recency-based eviction as a scoring-independent control (\S\ref{sec:scoring-control}) provides a principled falsification test that, to our knowledge, has no precedent in the KV cache compression literature. Prior work employs recency eviction (StreamingLLM~\cite{streamingllm-iclr2024}) or H$_2$O scoring~\cite{h2o-neurips2023} as performance baselines, but not as controls for isolating the contribution of the scoring mechanism itself. The null result ($d_z{=}0.083$, NS at $s{=}0.15$, $n{=}100$) confirms that the phase gradient is attention-score-specific rather than an artifact of phase segmentation or positional structure, strengthening the causal interpretation of strategy specificity as a property of score-dependent token selection.

\paragraph{Evidential status: mechanism discrimination, not causal proof.}
The triangulation across importance, random, and recency eviction constitutes \emph{mechanism discrimination}: we identify attention-score-based selection as the specific mediating pathway for phase-dependent severity, by showing the gradient's presence under score-based eviction and absence under score-independent controls. This is discriminative evidence, not intervention-based causal proof. Full causal demonstration would require directly manipulating score quality (e.g., injecting calibrated noise into attention scores before eviction) and observing a corresponding change in gradient magnitude, which we leave to future work. Our contribution is identifying the mediating pathway with sufficient specificity to guide system design: practitioners can use the finding that score-independent eviction eliminates the phase gradient to inform compression strategy selection, regardless of whether the underlying mechanism is fully causally resolved.
